\documentclass[11pt]{article}
\usepackage[margin=1in]{geometry}
\usepackage{amsmath,amssymb,amsthm,booktabs}
\usepackage[hidelinks]{hyperref}
\newtheorem{theorem}{Theorem}
\newcommand{\R}{\mathbb R}
\newcommand{\dotp}[2]{\langle #1,#2\rangle}
\title{A nonconvex incenter locus with nested circular caustics\\
\large A counterexample to MathDB 351604}
\author{Principia Math in Progress}
\date{12 September 2026}
\begin{document}
\maketitle
\section*{Task 1. Resolve the bic-III convexity conjecture}
Conjecture 2 of \cite[p.~23]{source} asserts that the incenter locus
of a bic-III triangle family is a convex curve. Its four circles must
be distinct. We use \emph{convex curve} in the usual boundary-of-a-convex-set
sense, rather than requiring that the curve itself be a convex set.

\begin{theorem}
There exists a bic-III family with all three circular caustics strictly
inside its circumcircle, and with no degenerate triangles, whose
incenter locus is not a convex curve.
\end{theorem}
\begin{proof}
Put $h=2/5$, $r=7/12$, and $A=1+h^2-r^2=2951/3600$.
Consider the pencil
\[
 C_\lambda:\quad x^2+y^2-2\lambda h x-1+\lambda A=0.
\]
Its center is $(\lambda h,0)$ and its squared radius is
$1-\lambda A+\lambda^2h^2$.
The circumcircle is $C_0$, the unit circle.
Use $C_1$ and $C_{7/4}$ as the first two caustics. Their centers and
radii are
\[
 O_1=(2/5,0),\quad r_1=7/12,\qquad
 O_2=(7/10,0),\quad r_2=\sqrt{799}/120.
\]
They are strictly nested:
$2/5+7/12<1$ and $3/10+\sqrt{799}/120<7/12$.
Their centers are distinct.

Here is a completely specified family, including the tangent choices.
Set
\[
 P(t)=\left(\frac{1-t^2}{1+t^2},\frac{2t}{1+t^2}\right),
 \qquad J(x,y)=(-y,x),\qquad t\in\R\cup\{\infty\}.
\]
For a circle with center $O$, radius $\rho$, and sign
$s\in\{-1,1\}$, define
\begin{align}
 V&=O-P,\qquad D=\dotp{V}{V},\qquad
 W=\frac{\sqrt{D-\rho^2}\,V+s\rho JV}{D}, \label{eq:W}\\
 T(O,\rho,s;t)&=P-2\dotp{P}{W}W. \label{eq:T}
\end{align}
Take $Q(t)=T(O_1,r_1,-1;t)$ and
$Z(t)=T(O_2,r_2,1;t)$.
The vector $W$ has length one, and the distance from $O$ to the
line $P+\R W$ is $\rho$. Equation \eqref{eq:T} gives its second
intersection with the unit circle. Thus $PQ$ and $PZ$ are the
specified tangent chords. The square roots are positive everywhere
because the caustics lie strictly inside the unit circle.
Also $P,Q,Z$ are pairwise distinct: $Q=Z$ would give a common
tangent to the two strictly nested caustics. Three distinct points
on a circle are noncollinear. Consequently the family and its
incenter are smooth, including through the parameter value at infinity.

Poncelet's Main Lemma, \cite[Lemma 1]{source}, gives a fixed
in-pencil envelope for $QZ$. We verify explicitly that it is a
third proper inner circle, distinct from the first two.
For either $t=0$ or $t=1$, let
\[
 n=J(Z-Q),\qquad e=\dotp{n}{Q}.
\]
The line $QZ$ is tangent to $C_\lambda$ precisely when
\begin{equation}
 g_t(\lambda):=(e-\lambda h n_x)^2-
 (1-\lambda A+\lambda^2h^2)\dotp{n}{n}=0. \label{eq:g}
\end{equation}
Exact rational interval evaluation gives the following isolating
intervals for the two roots; all endpoints are divided by $10^9$:
\[
\begin{array}{c|cc}
t&\text{smaller root}&\text{larger root}\\ \hline
0&(1957015893,1957015894)&(11272411549,11272411550)\\
1&(1957015893,1957015894)&(3366427888,3366427889)
\end{array}
\]
The coefficients of $\lambda^2$ are nonzero. Since the larger-root
intervals are disjoint, the fixed envelope supplied by the lemma must
be $C_{\lambda_*}$, where $\lambda_*$ is the smaller root of $g_0$.
In particular $1.95<\lambda_*<1.96$. Throughout that interval the
squared radius is between $0.001744444$ and $0.016197667$.
Writing its radius as $r_*$, the same interval evaluation gives
\[
 h\lambda_*+r_*<0.912<1,\qquad
 h\lambda_*-7/10+r_*<0.212<r_2.
\]
Thus the third caustic is strictly inside $C_{7/4}$.
Its parameter differs from $0,1,7/4$, so all four circles are distinct.

The incenter is, using ordinary positive side lengths,
\begin{equation}
 I(t)=\frac{aP+bQ+cZ}{a+b+c},\quad
 a=|Q-Z|,\quad b=|P-Z|,\quad c=|P-Q|. \label{eq:I}
\end{equation}
For vectors $u,v$, write $[u,v]=u_xv_y-u_yv_x$.
The following exact enclosures suffice for the counterexample:
\[
\begin{array}{c|c}
\text{quantity}&\text{enclosure}\\ \hline
I_x'(-1/3)&(206784605,206784606)/10^9\\
{}[I'(-1/3),\,I(0)-I(-1/3)]&(14992242,14992243)/10^9\\
{}[I'(-1/3),\,I(-1/2)-I(-1/3)]&(-144056,-144055)/10^9
\end{array}
\]
The first row shows that the tangent is nonzero. The next two rows
put points of the same incenter locus strictly on opposite sides of
that tangent line. If this locus were the boundary of a convex set,
a supporting line at $I(-1/3)$ would have to be its tangent: apply the
support inequality to $I(-1/3+\varepsilon)$ for both signs of
$\varepsilon$ and differentiate at zero. All boundary points would
then lie in one closed half-plane, contradicting the two signs.

For completeness, these are rigorous integer interval computations,
not floating-point sign tests. The accompanying \texttt{verify.py}
evaluates \eqref{eq:W}--\eqref{eq:I} with outward rounding to denominator
$B=10^{60}$ after every operation. An interval is stored as integer
endpoints $[l,u]/B$. Multiplication rounds the minimum and maximum
endpoint products outwards after division by $B$; division does the
same with the four endpoint quotients, rejecting a denominator
containing zero. Positive square roots are enclosed by
\[
 \left[\frac{\lfloor\sqrt{lB}\rfloor}{B},
 \frac{\lceil\sqrt{uB}\rceil}{B}\right].
\]
The integer square-root bounds are checked by squaring.
Derivatives are evaluated with pairs $(f,f')$ and the product,
quotient and square-root differentiation rules, starting with $(t,1)$.
The quadratic formula applied to \eqref{eq:g} gives the displayed
root intervals. These operations, using only integer arithmetic and
rational quotients, reproduce every strict bound above.
\end{proof}

The related paper \cite{hyperbolic} discusses cusps caused by exterior
caustics and degenerate triangles. The present counterexample uses
strictly nested proper circles and a smooth nondegenerate family;
it does not rely on those phenomena.

\section*{Task 2. Determine the convex parameter region}
Characterize the pairs of nested caustics in a circular pencil for
which a selected bic-III incenter locus is convex, taking account
of all four choices of tangent chords.

\section*{Task 3. Describe the onset of nonconvexity}
Find exact parameter conditions for the first vanishing of the signed
curvature, and determine which bifurcations preserve regularity and
simplicity of the incenter locus.

\section*{Task 4. Determine the algebraic locus}
Compute a minimal implicit equation and its degree for the incenter
locus of this example, separating the real incenter branch from
the excenter branches introduced by squaring.

\section*{Task 5. Compare the four triangle families}
Determine how convexity, inflections, and winding numbers change
under the four tangent choices associated with fixed first and
second caustics.

\begin{thebibliography}{9}
\bibitem{source}
R. Garcia, B. Odehnal and D. Reznik,
\emph{Loci of Poncelet triangles in the general closure case},
J. Geom. \textbf{113} (2022), article 17, 30 pp.,
\href{https://doi.org/10.1007/s00022-022-00629-3}{doi:10.1007/s00022-022-00629-3}.
Conjecture 2, p.~23; arXiv:2108.05430v2, p.~16.
\bibitem{hyperbolic}
R. Alves Garcia, B. Odehnal and D. Reznik,
\emph{Poncelet porisms in hyperbolic pencils of circles},
J. Geom. Graphics \textbf{25} (2021), 205--225.
\end{thebibliography}
\end{document}
